%==================== RESUMO ====================
\section*{RESUMO}
\addcontentsline{toc}{section}{RESUMO}

A elevada penetração da geração fotovoltaica na matriz elétrica brasileira, intensificada no Nordeste, transforma a variabilidade dessa fonte em um desafio operacional concreto, evidenciado pelos cortes de geração (\textit{curtailment}). Este trabalho propõe um modelo de previsão \textit{day-ahead} da geração de usinas fotovoltaicas baseado em aprendizado de máquina, sob abordagem puramente endógena — sem previsões numéricas do tempo, valendo-se apenas de variáveis observadas na usina e de âncoras determinísticas (posição solar e calendário). Utilizam-se dados reais do SCADA de uma usina de aproximadamente 2.5\,MW no Nordeste, com cerca de seis meses de cobertura. Comparam-se quatro modelos — XGBoost, SVM, MLP e LSTM — avaliados por RMSE, MAE e \textit{forecast skill score}. Este relatório consolida a revisão de literatura, o pré-processamento, a engenharia de \textit{features} e o protocolo experimental, apresentando o cronograma das etapas restantes.

\vspace{0.5cm}
\noindent\textbf{Palavras-chave:} energia fotovoltaica; aprendizado de máquina; previsão de geração; séries temporais; usinas solares.